\chapter{Extensions and Tasks}
\label{ch:tasks}

Some work takes longer than anybody wants to hold a connection open for.
Underwriting a loan, running a batch, waiting on a human being to make up their
mind. The Tasks extension is how a server says ``this will take a while, here is
a handle, come back for it'', and the second half of this chapter is how that
works in detail.

The first half is about the fact that Tasks arrives as an \emph{extension} at
all. The extensions framework is how MCP now grows without breaking
implementations that already exist, and Tasks is the worked example: it used to
be in core, and this revision moved it out.

\section{The extensions framework}

The problem every successful protocol hits: a feature that half the ecosystem
needs and the other half never will. Put it in core and you burden everyone. Put
it nowhere and it gets reinvented five incompatible ways.

MCP's answer is negotiated extensions, and the mechanism is deliberately small.
Both sides advertise support in a map keyed by identifier:

\begin{wire}
{
  "capabilities": {
    "tools": {},
    "extensions": {
      "io.modelcontextprotocol/tasks": {},
      "io.modelcontextprotocol/ui": {
        "mimeTypes": ["text/html;profile=mcp-app"]
      }
    }
  }
}
\end{wire}

Clients declare theirs per request, inside
\texttt{io.modelcontextprotocol/clientCapabilities}. Servers declare theirs in
\texttt{server/discover}. Identifiers follow the \texttt{\_meta} key rules, so
they carry a mandatory reverse-DNS prefix, and the
\texttt{io.modelcontextprotocol/} prefix is reserved for official ones.

The settings object is the extension's own business. Empty means ``supported,
nothing to configure''.

\subsection{The fallback rule}

One line, and it is the important one:

\keyidea{If one party supports an extension and the other does not, the
supporting party must either fall back to core behaviour or reject the request
with a clear error. It must never send a shape the other side cannot parse.}

Meridian's risk server does exactly this. It supports Tasks. It checks whether
the caller does, per request, and takes the synchronous path when they do not:

\begin{py}
def underwrite_loan(ctx: RequestContext):
    account_id = ctx.arguments["accountId"]
    account = ACCOUNTS.get(account_id)
    if account is None:
        return tool_error(f"No account {account_id}.")

    # Only return a task to a client that asked for the extension. To one
    # that did not, the fast path is the only correct answer.
    if not ctx.capabilities.supports_extension(tasks_ext.EXTENSION_ID):
        return _underwrite_now(account, ctx.arguments)

    task = server.tasks.create(status=tasks_ext.WORKING,
                               status_message="queued", poll_interval_ms=200)
    ...
\end{py}

Both paths are tested, because ``we fall back gracefully'' is a claim that rots:

\begin{py}
def test_client_without_the_extension_gets_the_synchronous_answer(self):
    server = risk.build_server()
    tasks_ext.install(server)
    client = Client(InProcessTransport(server),
                    capabilities=ClientCapabilities())
    result = client.call_tool("underwrite_loan",
                              {"accountId": "ACC-1000", "amountUsd": 250_000})
    self.assertEqual(result["resultType"], "complete")
    self.assertIn("decision", result["structuredContent"])
\end{py}

\subsection{Why this design is good}

Three reasons, and the third is the one that will matter in 2028.

\textbf{The core stays small, and features get somewhere to grow up.} Tasks
moved \emph{out} of core in this revision, which is a protocol getting smaller
and almost never happens. Read SEP-2663 for the reason, though, because it is
not that Tasks were judged unimportant. The opposite: the SEP calls them a
foundational building block and says outright that once the extension stabilises
and gets broad adoption, the intention is to promote it back into core.

The experimental \texttt{tasks} in \texttt{2025-11-25} was a stopgap shipped
before the extension mechanism existed. Moving it out buys it the ability to
evolve on its own cadence, on real implementation feedback, without being pinned
to the core specification's release schedule. That is what an extension is for,
and it is a more interesting claim than ``the core got smaller''.

\textbf{Extensions version independently.} MCP Apps was finalised
\texttt{2026-01-26} while core was on a different cycle. Neither had to wait.

\textbf{Vendors get a legitimate path.} Your company can ship
\texttt{com.acme/audit} without forking the protocol or squatting on a reserved
prefix. If it turns out to be generally useful, the road to standardisation runs through a public specification proposal.

\section{Why not just block?}

Now Tasks proper. And the fair question first: why not hold the connection open?

Sometimes you should. If the work takes 200 milliseconds, blocking is simpler
and cheaper and you should do that. Tasks solve four problems that blocking
cannot, and if none of them apply to you, do not use Tasks.

\begin{description}[leftmargin=1.4em, style=nextline, itemsep=3pt]
\item[Timeouts you do not control] Every intermediary between you and the client
has an opinion about how long a request may take. Load balancers, proxies, CDNs,
and mobile radios all reap idle connections, typically somewhere between 30 and
120 seconds.
\item[Crash resilience] A task id is durable. Client restarts, laptop lid
closes, network changes from wi-fi to cellular, and the same id still resolves.
A blocked request survives none of that, and Chapter~\ref{ch:transport} removed
stream resumability, so a broken stream loses the whole request.
\item[Progress visibility] A blocked request is opaque. A task has a status, a
message, and a progress fraction.
\item[Mid-flight input] The killer feature. A task can move to
\texttt{input\_required}, surface a question, and wait, without anyone holding a
connection open or inventing a second channel.
\end{description}

\section{The lifecycle}

\bookfig[0.86]{task-lifecycle}{Five states, three of them terminal. The two
notes at the bottom are the traps, and both of them have bitten
somebody.}{task-lifecycle}

\begin{table}[htbp]
\centering
\small
\begin{tabularx}{\textwidth}{@{}llX@{}}
\toprule
\thead{Status} & \thead{Terminal} & \thead{Meaning} \\
\midrule
\texttt{working}        & no  & In progress \\
\texttt{input\_required} & \textbf{no} & Needs client input. See \texttt{inputRequests} \\
\texttt{completed}      & yes & \texttt{result} holds what the call would have returned \\
\texttt{failed}         & yes & \texttt{error} holds a JSON-RPC error \\
\texttt{cancelled}      & yes & Cancelled, and cancellation is cooperative \\
\bottomrule
\end{tabularx}
\caption{The bolded row is the trap: a poller that stops when the status is not
\texttt{working} will hang forever on a task that is waiting for a human.}
\label{tab:task-states}
\end{table}

\subsection{Creation}

The server decides, per request, whether to return a task. There is no per-tool
flag and no per-request opt-in beyond the client's capability declaration. That
was a deliberate simplification in the redesign: the server knows how long the
work will take, and the client should not have to guess.

\begin{wire}
{
  "jsonrpc": "2.0",
  "id": 1,
  "result": {
    "resultType": "task",
    "task": {
      "taskId": "tsk_a1b2c3d4e5f6",
      "status": "working",
      "statusMessage": "queued",
      "ttlMs": 3600000,
      "pollIntervalMs": 200
    }
  }
}
\end{wire}

One requirement that sounds like an implementation detail and is not: the task
must be \textbf{durably created before the response is sent}. If the server
crashes between sending the id and persisting the task, the client holds a
handle to nothing, and it will poll that handle politely until its timeout.

\subsection{Polling}

\begin{wire}
{ "jsonrpc": "2.0", "id": 2, "method": "tasks/get",
  "params": { "taskId": "tsk_a1b2c3d4e5f6" } }
\end{wire}

Two things clients get wrong here, and Meridian's loop guards both:

\begin{py}
def poll_until_done(client, task_id, *, timeout=30.0,
                    on_status=None, input_provider=None):
    """Client-side polling loop, honouring the server's suggested cadence.

    Two things people get wrong here. First, ignoring `pollIntervalMs` and
    hammering every 50 ms, which turns one long request into four hundred short
    ones. Second, forgetting that `input_required` is not terminal: a task that
    stops to ask a question sits there until somebody answers it.
    """
    deadline = time.time() + timeout
    while time.time() < deadline:
        task = client.call("tasks/get", {"taskId": task_id}, use_cache=False)
        state = task.get("task", task)
        status = state.get("status")

        if status in TERMINAL:
            if status == FAILED:
                err = state.get("error") or {}
                raise errors.McpError(err.get("code", errors.INTERNAL_ERROR),
                                      err.get("message", "task failed"))
            return state.get("result") or {}

        if status == INPUT_REQUIRED and input_provider is not None:
            answers = {}
            for key, req in (state.get("inputRequests") or {}).items():
                if req.get("method") == "elicitation/create":
                    answers[key] = input_provider.elicit(key,
                                                         req.get("params") or {})
            client.call("tasks/update",
                        {"taskId": task_id, "inputResponses": answers},
                        use_cache=False)

        interval = max(0.05, state.get("pollIntervalMs", 500) / 1000.0)
        time.sleep(interval)
\end{py}

\begin{measurebox}{Polling cadence, for a 60-second job (arithmetic, not measured)}
\begin{tabular}{@{}rrrr@{}}
\toprule
\thead{Interval} & \thead{Polls} & \thead{Wasted requests} & \thead{Detection lag} \\
\midrule
50\,ms   & 1{,}200 & 1{,}199 & 50\,ms \\
500\,ms  & 120     & 119     & 500\,ms \\
2\,s     & 30      & 29      & 2\,s \\
10\,s    & 6       & 5       & 10\,s \\
\bottomrule
\end{tabular}

\vspace{6pt}
\footnotesize
Straight division, included because the shape of the trade is what matters. The
server picks \texttt{pollIntervalMs} because only the server knows the expected
duration.
\end{measurebox}

For long jobs, back off. Poll fast for the first few seconds in case the work
was quicker than expected, then widen. And add jitter, or a thousand clients
that started together will keep polling together forever.

\subsection{Or stop polling}

Every row in that table is waste, and the protocol has an answer.

A server may push task state with \texttt{notifications/tasks}, and clients opt
in through the \texttt{subscriptions/listen} mechanism from
Chapter~\ref{ch:resources}. Polling remains the default and the fallback,
because it needs nothing from either side. Push is what you use when the server
offers it.

The notification carries \textbf{the whole task}, not just its id:

\begin{wire}
{
  "jsonrpc": "2.0",
  "method": "notifications/tasks",
  "params": {
    "_meta": { "io.modelcontextprotocol/subscriptionId": 7 },
    "task": {
      "taskId": "tsk_a1b2c3d4e5f6",
      "status": "working",
      "statusMessage": "checking covenants",
      "progress": 0.75,
      "pollIntervalMs": 200
    }
  }
}
\end{wire}

That is the design decision worth noticing. A notification saying only ``task
tsk\_a1b2 changed'' would force a \texttt{tasks/get} to find out what changed,
so every update would cost a round trip and you would have replaced polling with
polling plus a wake-up. Carrying the state means the client learns everything
from one message it did not have to ask for.

Meridian pushes from inside \texttt{TaskStore.update}, which is the only path
that changes a task:

\begin{py}
def update(self, task_id: str, **fields) -> Task:
    """Change a task's state, and tell anybody who asked to be told.

    Every state change goes through here, which is what makes the push
    path trustworthy: there is no way to move a task to `completed` without
    the notification going out.
    """
\end{py}

The subscription filter has the same shape as the rest, and the granting logic
has one rule specific to this:

\begin{py}
# Task pushes are only on offer from a server running the Tasks
# extension. Granting them otherwise promises updates that no code
# path can ever send.
"tasks": "io.modelcontextprotocol/tasks" in (caps.get("extensions") or {}),
\end{py}

A client that asks a task-less server for task updates is told so in the
acknowledgement, immediately, rather than waiting quietly for the rest of its
life. That is the same courtesy the acknowledgement pays for resource
subscriptions, and for the same reason.

\begin{notebox}{Push does not let you skip the polling code}
Write the polling loop anyway.

A subscription is a stream, and streams break. The client reconnects and
re-sends \texttt{subscriptions/listen}, and in the window between the break and
the re-subscribe, any number of task transitions happened that nobody delivered.
The state you missed is still readable with \texttt{tasks/get}, which is why the
right shape is push as an optimisation over a polling loop with a long interval,
rather than push as a replacement for one.
\end{notebox}

\subsection{Mid-flight input}

This is where Tasks and MRTR compose, and it is genuinely elegant.

A task hits a point needing human judgement. It moves to
\texttt{input\_required} and surfaces the same \texttt{inputRequests} map from
Chapter~\ref{ch:mrtr}. The client answers with \texttt{tasks/update} rather than
by retrying the original call, because the original call is long since finished.

\begin{wire}
{
  "jsonrpc": "2.0",
  "id": 5,
  "method": "tasks/update",
  "params": {
    "taskId": "tsk_a1b2c3d4e5f6",
    "inputResponses": {
      "approval": { "action": "accept",
                    "content": { "approver": "j.okonjo" } }
    }
  }
}
\end{wire}

The server acknowledges with an empty result, and one detail is worth copying:

\begin{py}
@server.method("tasks/update")
def _update(ctx: RequestContext) -> dict:
    """Deliver answers to a task waiting on input.

    Unknown or already-satisfied keys are ignored rather than rejected. A
    client that retried a `tasks/update` after a timeout should not get an
    error for being careful.
    """
\end{py}

Idempotent updates. A client whose \texttt{tasks/update} timed out will retry,
and punishing it for that turns a transient network problem into a failed task.

\subsection{Cancellation is a request, not a command}

\begin{wire}
{ "jsonrpc": "2.0", "id": 6, "method": "tasks/cancel",
  "params": { "taskId": "tsk_a1b2c3d4e5f6" } }
\end{wire}

The server acknowledges the \emph{intent}. It is not obliged to stop, and the
task may still reach \texttt{completed}. That is not sloppiness: a task that has
already charged a card cannot be un-charged by a cancel arriving 200
milliseconds later, and pretending otherwise would be worse.

Meridian's test asserts the honest condition:

\begin{py}
def test_cancel_is_acknowledged(self):
    ...
    client.call("tasks/cancel", {"taskId": task_id}, use_cache=False)
    self.assertIn(store.get(task_id).status,
                  (tasks_ext.CANCELLED, tasks_ext.COMPLETED))
\end{py}

Either outcome is correct. A test that demanded \texttt{CANCELLED} would be
asserting something the protocol does not promise.

\section{Where tasks live, and why it matters}

A task outlives the request that created it, so something has to store it. Where
that something lives decides whether your deployment survives having more than
one replica.

\begin{dangerbox}{An in-memory task store reinvents sticky sessions}
The client polls with \texttt{tasks/get}. Behind a load balancer, that poll
lands on whichever replica the balancer chooses, which is almost certainly not
the one that created the task.

If your task store is a process-local dictionary, the poll returns ``unknown
task'' and the work is lost. The fix is not session affinity. The fix is a
shared, durable store, because the whole point of statelessness was to stop
needing affinity.
\end{dangerbox}

Meridian is honest about being a book:

\begin{py}
class TaskStore:
    """Where tasks live between polls.

    In-memory here, which is a lie a book is allowed to tell once. Any real
    deployment needs this in something durable and shared, because the whole
    point is that the next poll may land on a different replica. The moment you
    put tasks in a process-local dict behind a load balancer, you have
    reinvented sticky sessions, which is the thing this revision deleted.
    """
\end{py}

Requirements for a real one:

\begin{itemize}[leftmargin=1.6em, itemsep=3pt]
\item \textbf{Shared} across replicas. Redis, Postgres, DynamoDB, anything.
\item \textbf{Durable} enough to survive a deploy. Task ids outlive processes,
which is the feature.
\item \textbf{Expiring.} \texttt{ttlMs} is a promise about how long the id
resolves. Sweep afterwards or the store grows forever.
\item \textbf{Authorized on every access.} A task id is a name, not a
capability. Check the caller on \texttt{tasks/get} as well as on creation, or
anyone who can guess an id can read somebody's underwriting decision.
\end{itemize}

\section{When not to use Tasks}

Tasks add a polling loop, a durable store, an expiry policy, and a second code
path in every client. That is real complexity, and for short work it buys
nothing.

\begin{table}[htbp]
\centering
\small
\begin{tabularx}{\textwidth}{@{}lXX@{}}
\toprule
\thead{Duration} & \thead{Use} & \thead{Because} \\
\midrule
Under 1\,s      & Synchronous & Polling costs more than the work \\
1 to 10\,s      & Synchronous, with progress notifications & The stream is still healthy; the user sees movement \\
10 to 60\,s     & Judgement call & Depends on your intermediaries' patience \\
Over 60\,s      & Tasks & Something in the path will reap the connection \\
Needs a human   & Tasks & Humans take minutes, and connections do not \\
Wraps a job API & Tasks & You already have a job id; map it \\
\bottomrule
\end{tabularx}
\caption{Thresholds. The 10-to-60-second band is where measurement beats
opinion.}
\label{tab:when-tasks}
\end{table}

For the middle band, use progress notifications instead. They keep the user
informed without any of the machinery:

\begin{py}
scored, unknown = [], []
total = len(ids)
for i, account_id in enumerate(ids):
    ...
    if total > 20 and i % 10 == 0:
        ctx.progress(i, total, f"scored {i} of {total}")
\end{py}

Note the guard. Progress on a three-item loop is noise, and each notification is
a frame on the wire that somebody has to read.

\begin{notebox}{Progress and tasks are different channels}
\texttt{notifications/progress} flows on the \emph{response stream of the
request it belongs to}. It never appears on a \texttt{subscriptions/listen}
stream, and it stops existing the moment the request finishes.

Task status is polled, or pushed via \texttt{notifications/tasks} to clients
who subscribed, and it outlives the request that created it by design.
Different lifetime, different channel, and mixing them up produces a client
that waits for progress on a request that ended twenty minutes ago.
\end{notebox}

\section{Meridian's underwriting task}

Four stages, progress reported at each:

\begin{py}
task = server.tasks.create(status=tasks_ext.WORKING,
                           status_message="queued", poll_interval_ms=200)

def work(t):
    stages = ["pulling filings", "scoring", "checking covenants", "pricing"]
    for i, stage in enumerate(stages, start=1):
        server.tasks.update(t.task_id, status_message=stage,
                            progress=i / len(stages))
        time.sleep(0.05)
    return _underwrite_now(account, ctx.arguments)

tasks_ext.run_in_background(server.tasks, task, work)
return tasks_ext.create_task_result(task)
\end{py}

And the background runner, whose only real job is making sure the task always
reaches a terminal state:

\begin{py}
def runner() -> None:
    try:
        result = work(task)
        if task.status not in TERMINAL:
            store.update(task.task_id, status=COMPLETED,
                         result=result if isinstance(result, dict) else {},
                         status_message="done")
    except errors.McpError as exc:
        store.update(task.task_id, status=FAILED, error=exc.to_json())
    except Exception as exc:
        store.update(task.task_id, status=FAILED,
                     error=errors.InternalError(str(exc)).to_json())
\end{py}

The bare \texttt{except} is deliberate. A task that crashes without reaching a
terminal state is a client polling forever, and forever is a long time.
\texttt{failed} with an explanation is always better than \texttt{working} with a
lie.

Note also the \texttt{if task.status not in TERMINAL} guard. If the work
completed after a cancellation was recorded, the cancellation stands. Terminal
means terminal.

\section{Vendor extensions}

You can define your own. Two rules and one piece of advice.

\textbf{Prefix it properly.} Reverse DNS, and the second label must not be
\texttt{modelcontextprotocol} or \texttt{mcp}. \texttt{com.acme/audit} is fine.
\texttt{com.mcp.acme/audit} is reserved and using it is squatting.

\textbf{Document the fallback.} What should a peer that does not support it do?
Answer that in your specification, not in a support ticket.

And the advice: before you write one, check whether a tool would do. An
extension changes the protocol for everyone who implements it. A tool is a tool.
The bar for a new extension should be that the capability genuinely cannot be
expressed as a tool, resource, or prompt, which is a higher bar than it first
appears.

\section{What to remember}

\begin{itemize}[leftmargin=1.6em, itemsep=3pt]
\item Extensions are negotiated through capabilities, with mandatory reverse-DNS
identifiers.
\item The fallback rule: revert to core behaviour or reject clearly. Never send
a shape the peer cannot parse, and test both paths.
\item Tasks moved out of core into \texttt{io.modelcontextprotocol/tasks}. A
protocol got smaller.
\item The server decides per request whether to return a task; the client opts
in once via capabilities.
\item Create the task durably \emph{before} answering.
\item \texttt{input\_required} is not terminal. Poll at the server's
\texttt{pollIntervalMs}, back off, and add jitter.
\item \texttt{tasks/update} must be idempotent. Cancellation is cooperative.
\item \texttt{notifications/tasks} carries the whole task, so a push costs no
follow-up round trip. Keep the polling loop anyway; streams break.
\item An in-memory task store behind a load balancer has reinvented sticky
sessions. Use a shared, durable, expiring store, and authorize every access.
\item Under a second, stay synchronous. One to ten seconds, stream progress.
Over a minute, or human in the loop, use Tasks.
\end{itemize}

Next: MCP Apps, where a tool result stops being text and starts being an
interface.
